\let\textcircled=\pgftextcircled
\chapter{Conception et mod\'elisation}
\label{chap:conception}

\textit{\initial{C}e chapitre constitue le c\oe ur du travail. Il expose les
principes qui gouvernent la conception, puis d\'eroule la mod\'elisation du syst\`eme
selon les trois vues du formalisme UML~2.5.1 : la vue statique, qui dit de quoi le
syst\`eme est fait ; la vue dynamique, qui dit comment il se comporte ; la vue
architecturale, qui dit comment il est d\'ecoup\'e et d\'eploy\'e. Il se cl\^ot sur la
conception de la persistance et sur la matrice reliant chaque exigence \`a
l'\'el\'ement de conception qui la porte.}

\mtcselectlanguage{french}
\minitoc
\newpage

%%%==========================================================================
\section{Principes directeurs}
%%%==========================================================================

\subsection{Une question de conception unique}
La conception r\'epond \`a une question, et une seule : \textbf{par quoi remplace-t-on
l'analyste qui relisait ?} Dans l'organisation ant\'erieure, une erreur du moteur
produisait une recommandation douteuse qu'un humain \'ecartait ; le co\^ut d'une
erreur \'etait le temps de lecture d'un analyste. D\`es lors que le syst\`eme ex\'ecute
lui-m\^eme, la m\^eme erreur produit une r\`egle de pare-feu qui coupe un service de
production. Ce n'est pas un changement de module : c'est un changement de nature.

Quatre m\'ecanismes se substituent \`a la relecture humaine. Ils ne sont pas
interchangeables : chacun couvre un mode de d\'efaillance distinct, et le retrait de
l'un ouvre une br\`eche que les autres ne comblent pas.

\begin{table}[H]
\centering
\caption{Modes de d\'efaillance et m\'ecanismes de conception associ\'es}
\label{tab:modes-defaillance}
\small
\begin{tabular}{@{}p{5.6cm}p{4.4cm}p{3.6cm}@{}}
\toprule
\textbf{Mode de d\'efaillance} & \textbf{M\'ecanisme} & \textbf{Exigence} \\
\midrule
Le syst\`eme agit sur une menace qu'il ne comprend pas & Garde d'ancrage documentaire & EF-04 \\
Le syst\`eme fait quelque chose d'irrattrapable & Catalogue de r\'eversibilit\'e & EF-14 \\
Le syst\`eme fait quelque chose que le CIRT ne veut pas & Politique compil\'ee a priori & EF-15 \\
Le syst\`eme a raison sur la menace mais nuit \`a la cible & Boucle de contr\^ole ferm\'ee & EF-25 \\
Le syst\`eme se trompe en boucle & Coupe-circuit global & EF-26 \\
Le syst\`eme se tait devant l'inconnu & Confinement par indicateurs & EF-27 \\
\bottomrule
\end{tabular}
\end{table}

Aucun de ces m\'ecanismes ne r\'eintroduit d'attente humaine. Ce sont des
\emph{refus} ou des \emph{corrections}, jamais des mises en attente : le mod\`ele
d'\'etats de l'incident (figure~\ref{fig:etats-incident}) ne comporte d\'elib\'er\'ement
aucun \'etat \og en attente de validation \fg{} sur le chemin nominal.

\subsection{Les cinq axes de conception}
\label{sec:axes}
Les d\'ecisions structurantes se rattachent \`a cinq axes, qui servent de crit\`ere
d'arbitrage chaque fois qu'un choix se pr\'esente.

\begin{table}[H]
\centering
\caption{Axes de conception et cons\'equence structurelle de chacun}
\label{tab:axes}
\small
\begin{tabular}{@{}p{0.8cm}p{4.4cm}p{8.4cm}@{}}
\toprule
& \textbf{Axe} & \textbf{Cons\'equence sur la structure} \\
\midrule
1 & Autonomie sous garanties & Aucun \'etat d'attente humaine sur le chemin nominal. \\
2 & \textbf{La r\'eversibilit\'e conditionne l'action} & La r\'eversibilit\'e est un attribut de \texttt{ActionSpec}, contr\^ol\'e \`a la construction, non une m\'etadonn\'ee documentaire. \\
3 & Explicabilit\'e opposable & Toute d\'ecision porte une \texttt{DecisionTrace} : sources, r\`egles appliqu\'ees, alternatives \'ecart\'ees. \\
4 & Priorisation par l'enjeu & L'incident calcule lui-m\^eme un score de risque d\'eterministe et explicable. \\
5 & Fonctionnement souverain & Aucune d\'ependance sortante obligatoire ; persistance locale ; mode d\'egrad\'e natif. \\
\bottomrule
\end{tabular}
\end{table}

\subsection{Style architectural}
Le style retenu est une \textbf{architecture hexagonale} (ports et adaptateurs).
Son int\'er\^et ici n'est pas th\'eorique : il tient \`a une propri\'et\'e unique --- le
\emph{paquetage du domaine ne d\'epend de rien}. Ni de la base de donn\'ees, ni du
cadriciel web, ni des connecteurs vers les \'equipements. Cette absence de
d\'ependance rend ses invariants incontournables : quel que soit le chemin d'appel
--- interface web, ligne de commande, test automatis\'e, rejeu de la file d\'egrad\'ee
--- une \texttt{ActionSpec} qui se dit r\'eversible sans porter de verbe d'annulation
ne peut pas exister.

\subsection{Patrons de conception mis en \oe uvre}
\begin{table}[H]
\centering
\caption{Patrons de conception et motif de leur emploi}
\label{tab:patrons}
\small
\begin{tabular}{@{}p{2.7cm}p{4.3cm}p{6.6cm}@{}}
\toprule
\textbf{Patron} & \textbf{Emploi} & \textbf{Motif du choix} \\
\midrule
Adaptateur & Normalisation des sources & Ajouter une source de d\'etection sans toucher au moteur (ENF-08). \\
Strat\'egie & Actionneurs & Chaque \'equipement expose le m\^eme contrat ; le moteur ignore la technologie sous-jacente. \\
Commande + M\'emento & Action et jeton d'annulation & L'annulation exige de conserver l'information n\'ecessaire au retour arri\`ere, produite \`a l'ex\'ecution. \\
D\'ep\^ot (\emph{Repository}) & Persistance & Isole le domaine du sch\'ema relationnel et rend le rejeu testable. \\
Coupe-circuit & Suspension de l'autonomie & Emp\^eche l'entretien d'une boucle d'erreurs (EF-26). \\
Cha\^ine de responsabilit\'e & Cha\^ine de d\'ecision & Chaque garde peut interrompre la cha\^ine en motivant son refus. \\
Objet-valeur & \texttt{Asset}, \texttt{DetectionEvent}, \texttt{ActionSpec} & Immuables : une observation ne se modifie pas, elle se remplace. \\
Fabrique & \texttt{Incident.from\_event()} & Garantit la coh\'erence de la cl\'e de corr\'elation d\`es la cr\'eation. \\
\bottomrule
\end{tabular}
\end{table}

%%%==========================================================================
\section{Vue statique}
%%%==========================================================================

\subsection{Diagramme de paquetages}
Le d\'ecoupage en paquetages mat\'erialise le style hexagonal. La lecture du
diagramme se fait par le sens des fl\`eches : \textbf{toutes convergent vers
\texttt{domain}}, aucune n'en sort. C'est la v\'erification la plus rapide de la
propri\'et\'e recherch\'ee, et elle est contr\^ol\'ee automatiquement au moment des tests.

\begin{figure}[H]
\centering
\begin{diagramme}[node distance=6mm]
  % --- couche interface
  \node (api)  [paquetage=3.1cm] at (0,0)   {\textbf{api}\\[1pt]\tiny routes REST, flux SSE, contr\^ole de r\^ole};
  \node (web)  [paquetage=3.1cm] at (4,0)   {\textbf{web}\\[1pt]\tiny interface d'exploitation};
  \node (cli)  [paquetage=3.1cm] at (8,0)   {\textbf{cli}\\[1pt]\tiny exploitation en ligne de commande};

  % --- couche application
  \node (orch) [paquetage=3.1cm] at (0,-2.6) {\textbf{orchestration}\\[1pt]\tiny d\'ecision, politique, repli, coupe-circuit};
  \node (asi)  [paquetage=3.1cm] at (4,-2.6) {\textbf{assistant}\\[1pt]\tiny faits, rapports, conversation};
  \node (det)  [paquetage=3.1cm] at (8,-2.6) {\textbf{detection}\\[1pt]\tiny UEBA, sondes, boucle post-action};

  % --- noyau
  \node (dom)  [paquetage=5.2cm, umlaccentue] at (4,-5.4)
        {\textbf{domain}\\[1pt]\tiny entit\'es, objets-valeurs, \'enum\'erations, invariants\\[1pt]\tiny\itshape ne d\'epend d'aucun autre paquetage};

  % --- couche infrastructure
  \node (ing)  [paquetage=3.1cm] at (-1.2,-8.2) {\textbf{ingestion}\\[1pt]\tiny adaptateurs de sources};
  \node (enr)  [paquetage=3.1cm] at (2.6,-8.2)  {\textbf{enrichment}\\[1pt]\tiny fonds documentaire, ancrage};
  \node (act)  [paquetage=3.1cm] at (6.4,-8.2)  {\textbf{actuators}\\[1pt]\tiny connecteurs d'\'equipements};
  \node (per)  [paquetage=3.1cm] at (10.2,-8.2) {\textbf{persistence}\\[1pt]\tiny d\'ep\^ots, sch\'ema};
  \node (aud)  [paquetage=3.1cm] at (10.2,-5.4) {\textbf{audit}\\[1pt]\tiny journal cha\^in\'e, notifications};

  \foreach \n in {api,web,cli,orch,asi,det,ing,enr,act,per,aud}{}

  \draw[dependance] (api) -- (orch);
  \draw[dependance] (web) -- (api);
  \draw[dependance] (cli) -- (orch);
  \draw[dependance] (asi) -- (dom);
  \draw[dependance] (det) -- (dom);
  \draw[dependance] (orch) -- (dom);
  \draw[dependance] (ing) -- (dom);
  \draw[dependance] (enr) -- (dom);
  \draw[dependance] (act) -- (dom);
  \draw[dependance] (per) -- (dom);
  \draw[dependance] (aud) -- (dom);
  \draw[dependance] (orch) -- (act);
  \draw[dependance] (orch) -- (aud);
\end{diagramme}
\caption{Diagramme de paquetages --- toutes les d\'ependances convergent vers le domaine}
\label{fig:paquetages}
\end{figure}

\subsection{Diagramme de classes du domaine}
\label{sec:classes-domaine}
La figure~\ref{fig:classes-domaine} pr\'esente les classes du noyau m\'etier. Trois
d\'ecisions de mod\'elisation y sont lisibles et m\'eritent d'\^etre justifi\'ees.

\paragraph{Immuabilit\'e des observations.} \texttt{Asset}, \texttt{DetectionEvent}
et \texttt{ActionSpec} sont des objets-valeurs immuables. Un \'ev\'enement de
d\'etection d\'ecrit un \emph{constat} \`a un instant donn\'e : le modifier apr\`es coup
reviendrait \`a r\'e\'ecrire l'histoire, ce qui ruinerait l'opposabilit\'e du journal
(ENF-02). \`A l'inverse, \texttt{Incident} et \texttt{ActionResult} sont mutables :
ce sont des \emph{\'etats} qui \'evoluent, et leurs transitions sont contraintes par
les machines \`a \'etats de la section~\ref{sec:etats}.

\paragraph{La r\'eversibilit\'e comme invariant de construction.} L'attribut
\texttt{reversibility} de \texttt{ActionSpec} n'est pas une annotation : le
constructeur refuse toute instance qui se d\'eclare r\'eversible sans porter de
verbe d'annulation. La contrainte est donc port\'ee par le type, non par le code
appelant --- aucun chemin d'ex\'ecution ne peut la contourner.

\paragraph{La trace n'est pas un journal.} \texttt{DecisionTrace} est un attribut
de la d\'ecision, pas une \'ecriture de journal. La distinction importe : le journal
enregistre \emph{qu'une} d\'ecision a eu lieu ; la trace porte \emph{pourquoi} elle a
\'et\'e prise --- sources documentaires, conditions satisfaites, actions envisag\'ees
puis \'ecart\'ees, verdicts de politique. C'est cette derni\`ere qui rend la d\'ecision
rejugeable.

\begin{sidewaysfigure}
\centering
\ajusterpaysage{%
\begin{diagramme}

  %%% ---------------------------------------------------- objets-valeurs
  \node (asset) [umlclasse=3.9cm] at (0,0) {
    \umltitres{value object}{Asset}
    \nodepart{two}
    - asset\_id \,: String\\
    - hostname \,: String [0..1]\\
    - ip \,: String [0..1]\\
    - user \,: String [0..1]\\
    - criticality \,: Integer \{1..5\}\\
    - zone \,: String
    \nodepart{three}
    + correlationKey() \,: String
  };

  \node (event) [umlclasse=4.6cm] at (5.6,0) {
    \umltitres{value object}{DetectionEvent}
    \nodepart{two}
    - eventId \,: String \{id\}\\
    - occurredAt \,: DateTime\\
    - receivedAt \,: DateTime\\
    - source \,: SourceKind\\
    - category \,: String\\
    - severity \,: Severity\\
    - confidence \,: Real \{0..1\}\\
    - indicators \,: Map\\
    - mitreTechniques \,: String [*]
    \nodepart{three}
    + fingerprint() \,: String\\
    + toDict() \,: Map
  };

  %%% ---------------------------------------------------------- agregat
  \node (incident) [umlclasse=4.8cm] at (11.9,0) {
    \umltitres{entity, aggregate root}{Incident}
    \nodepart{two}
    - incidentId \,: String \{id\}\\
    - correlationKey \,: String\\
    - status \,: IncidentStatus\\
    - severity \,: Severity\\
    - openedAt, updatedAt \,: DateTime\\
    - attackCode, family \,: String\\
    - dangerousness \,: Real\\
    - priorityRank \,: Integer
    \nodepart{three}
    + fromEvent(e) \,: Incident \guillemotleft factory\guillemotright\\
    + accepts(e) \,: Boolean\\
    + absorb(e)\\
    + riskScore() \,: Real
  };

  \node (attaque) [umlcompact=4.2cm] at (18.2,0) {
    \umltitres{reference data}{AttackType}
    \nodepart{two}
    - code \,: String \{A1..D4\}\\
    - family \,: AttackFamily\\
    - baseDangerousness \,: Real\\
    - reversibility \,: Reversibility\\
    - autonomouslyActionable \,: Boolean
  };

  %%% -------------------------------------------------------- decision
  \node (decision) [umlclasse=4.6cm] at (11.9,-8.4) {
    \umltitres{entity}{Decision}
    \nodepart{two}
    - decisionId \,: String \{id\}\\
    - incidentId \,: String\\
    - outcome \,: DecisionOutcome\\
    - rationale \,: String\\
    - decidedAt \,: DateTime\\
    - classification \,: Map\\
    - fallback \,: Map [0..1]
    \nodepart{three}
    + isActionable() \,: Boolean\\
    + toDict() \,: Map
  };

  \node (trace) [umlclasse=4.4cm] at (18.2,-8.4) {
    \umltitres{value object}{DecisionTrace}
    \nodepart{two}
    - playbookId, version \,: String\\
    - matchedConditions \,: String [*]\\
    - groundingScore \,: Real\\
    - contextSources \,: String [*]\\
    - consideredActions \,: String [*]\\
    - rejectedActions \,: Map [*]
    \nodepart{three}
    + toDict() \,: Map
  };

  \node (verdict) [umlcompact=3.6cm] at (18.2,-13.6) {
    \umltitres{value object}{PolicyVerdict}
    \nodepart{two}
    - allowed \,: Boolean\\
    - ruleId \,: String [0..1]\\
    - ruleText \,: String [0..1]\\
    - reason \,: String
  };

  %%% ---------------------------------------------------------- actions
  \node (spec) [umlclasse=4.6cm] at (5.6,-8.4) {
    \umltitres{value object}{ActionSpec}
    \nodepart{two}
    - verb \,: String\\
    - actuator \,: String\\
    - target \,: String\\
    - parameters \,: Map\\
    - reversibility \,: Reversibility\\
    - rollbackVerb \,: String [0..1]\\
    - blastRadius \,: Integer \{$\geq$1\}\\
    - expectedEffect \,: String
    \nodepart{three}
    + key() \,: String
  };

  \node (result) [umlclasse=4.4cm] at (0,-8.4) {
    \umltitres{entity}{ActionResult}
    \nodepart{two}
    - actionId \,: String \{id\}\\
    - status \,: ActionStatus\\
    - startedAt, finishedAt \,: DateTime\\
    - rollbackToken \,: String [0..1]\\
    - rolledBackAt \,: DateTime [0..1]\\
    - rollbackActor \,: String\\
    - error \,: String [0..1]
    \nodepart{three}
    + durationMs() \,: Integer\\
    + isReversible() \,: Boolean\\
    + markStarted() / markExecuted()
  };

  %%% ------------------------------------------------------- contraintes
  \node (inv1) [note=5.2cm] at (5.6,-14.2) {
    \textbf{\{invariant de construction\}}\\
    \texttt{reversibility} $\neq$ \textsc{irreversible}\\
    $\Rightarrow$ \texttt{rollbackVerb} $\neq$ \textit{null}\\[2pt]
    Sans verbe d'annulation, une action d\'eclar\'ee
    r\'eversible rendrait la boucle de contr\^ole (EF-25)
    inop\'erante. L'instance est refus\'ee \`a la construction.
  };
  \node (inv2) [note=4.6cm] at (0,-14.2) {
    \textbf{\{invariant\}}\\
    \texttt{isReversible()} $\equiv$
    \texttt{spec.reversibility} $\neq$ \textsc{irreversible}
    $\wedge$ \texttt{rollbackToken} $\neq$ \textit{null}\\[2pt]
    Le jeton n'existe qu'apr\`es ex\'ecution r\'eussie.
  };

  %%% ------------------------------------------------------- relations
  \draw[composition] (event.west) -- node[mult, near start]{1} node[mult, near end]{1} (asset.east);
  \draw[composition] (incident.west) -- node[mult, near start]{1} node[mult, near end]{1..*} (event.east);
  \draw[dependance]  (incident.east) -- node[role]{classifie par} (attaque.west);
  \draw[composition] (incident.south west) -- ++(0,-1.2) -| node[mult, near start]{1} node[mult, near end]{0..*} (result.north);
  \draw[association] (decision.north) -- node[mult, near start]{0..*} node[mult, near end]{1} (incident.south);
  \draw[composition] (decision.east) -- node[mult, near start]{1} node[mult, near end]{1} (trace.west);
  \draw[composition] (trace.south) -- node[mult, near start]{1} node[mult, near end]{0..*} (verdict.north);
  \draw[agregation]  (decision.west) -- node[mult, near start]{1} node[mult, near end]{0..*} (spec.east);
  \draw[association] (result.east) -- node[mult, near start]{0..*} node[mult, near end]{1} (spec.west);
  \draw[dependance]  (inv1.north) -- (spec.south);
  \draw[dependance]  (inv2.north) -- (result.south);
\end{diagramme}}
\caption{Diagramme de classes du domaine}
\label{fig:classes-domaine}
\end{sidewaysfigure}

\subsection{\'Enum\'erations du domaine}
Les \'enum\'erations forment le vocabulaire ferm\'e du syst\`eme. Leurs valeurs sont
s\'erialis\'ees telles quelles dans le journal d'audit : elles font partie du contrat
de tra\c cabilit\'e et ne peuvent \^etre renomm\'ees sans une entr\'ee d'historique de
version. La figure~\ref{fig:enumerations} les recense.

\begin{figure}[H]
\centering
\begin{diagramme}
  \node (sev) [umlcompact=3.2cm] at (0,0) {
    \umltitres{enumeration}{Severity}
    \nodepart{two}
    INFO\\ LOW\\ MEDIUM\\ HIGH\\ CRITICAL
  };
  \node (rev) [umlcompact=3.9cm] at (4.2,0) {
    \umltitres{enumeration}{Reversibility}
    \nodepart{two}
    REVERSIBLE\\ PARTIALLY\_REVERSIBLE\\ IRREVERSIBLE
  };
  \node (ist) [umlcompact=3.2cm] at (8.6,0) {
    \umltitres{enumeration}{IncidentStatus}
    \nodepart{two}
    OPEN\\ CONTAINED\\ ROLLED\_BACK\\ CLOSED
  };
  \node (ast) [umlcompact=3.6cm] at (12.4,0) {
    \umltitres{enumeration}{ActionStatus}
    \nodepart{two}
    PLANNED\\ BLOCKED\_BY\_POLICY\\ BLOCKED\_BY\_BREAKER\\ EXECUTING\\ EXECUTED\\ FAILED\\ ROLLED\_BACK\\ ROLLBACK\_FAILED
  };
  \node (dou) [umlcompact=4.4cm] at (2.1,-4.6) {
    \umltitres{enumeration}{DecisionOutcome}
    \nodepart{two}
    AUTONOMOUS\_EXECUTION\\ NO\_ACTION\_NEEDED\\ NO\_GROUNDED\_CONTEXT\\ POLICY\_DENIED\\ BREAKER\_OPEN\\ OUT\_OF\_CATALOG
  };
  \node (src) [umlcompact=3.2cm] at (6.8,-4.6) {
    \umltitres{enumeration}{SourceKind}
    \nodepart{two}
    UEBA\\ INFRASTRUCTURE\\ EDR\\ NIDS\\ SIEM\\ MANUAL\\ REPLAY
  };
  \node (fam) [umlcompact=3.6cm] at (11.2,-4.6) {
    \umltitres{enumeration}{AttackFamily}
    \nodepart{two}
    NETWORK \tiny(A)\normalsize\\ APPLICATION \tiny(B)\normalsize\\ INSIDER \tiny(C)\normalsize\\ INFRASTRUCTURE \tiny(D)\normalsize
  };
\end{diagramme}
\caption{\'Enum\'erations du domaine}
\label{fig:enumerations}
\end{figure}

\subsection{Contrats et invariants}
Le tableau~\ref{tab:invariants} \'enonce les contraintes que le mod\`ele garantit,
dans une forme proche d'OCL. Chacune est v\'erifi\'ee \`a la construction ou \`a la
transition, jamais laiss\'ee \`a la discipline de l'appelant.

\begin{table}[H]
\centering
\caption{Invariants du mod\`ele et niveau auquel ils sont garantis}
\label{tab:invariants}
\scriptsize
\begin{tabular}{@{}p{2.4cm}p{7.6cm}p{3.6cm}@{}}
\toprule
\textbf{Classe} & \textbf{Contrainte} & \textbf{Garanti \`a} \\
\midrule
\texttt{DetectionEvent} & \verb|confidence >= 0.0 and confidence <= 1.0| & la construction \\
\texttt{DetectionEvent} & \verb|asset.criticality in 1..5| & la construction \\
\texttt{ActionSpec} & \verb|reversibility <> IRREVERSIBLE implies| \newline \texttt{rollbackVerb <> null} & la construction \\
\texttt{ActionSpec} & \verb|blastRadius >= 1| & la construction \\
\texttt{ActionResult} & \verb|isReversible() implies rollbackToken <> null| & d\'eriv\'ee \\
\texttt{Incident} & \verb|accepts(e) implies| \newline \verb|status not in {CLOSED, ROLLED_BACK}| & la transition \\
\texttt{Incident} & \verb|riskScore() in 0..100| & le calcul \\
\texttt{Decision} & \verb|isActionable() implies| \newline \verb|outcome = AUTONOMOUS_EXECUTION and actions->notEmpty()| & d\'eriv\'ee \\
\textit{global} & \verb|Decision.actions->forAll(a |\verb#|# \verb|a.reversibility <> IRREVERSIBLE)| & la politique \\
\bottomrule
\end{tabular}
\end{table}

L'invariant global de la derni\`ere ligne est le plus important du mod\`ele :
\textbf{aucune action irr\'eversible ne peut figurer dans une d\'ecision ex\'ecutable}.
Il est v\'erifi\'e non pas sur un cas mais sur l'int\'egralit\'e du catalogue m\'etier, de
sorte qu'une entr\'ee ajout\'ee ult\'erieurement ne puisse franchir cette limite en
silence.

\subsection{Diagramme d'objets}
Un diagramme de classes dit ce qui est possible ; il ne dit pas ce qui se produit.
La figure~\ref{fig:objets} donne un instantan\'e du mod\`ele apr\`es traitement d'un
cas concret --- une exfiltration de donn\'ees d\'etect\'ee par une sonde r\'eseau sur un
serveur applicatif. On y lit l'effet des choix pr\'ec\'edents : deux actions ont \'et\'e
envisag\'ees, une seule ex\'ecut\'ee, et la trace conserve le motif du rejet de l'autre.

\begin{figure}[H]
\centering
\begin{diagramme}
  \node (o1) [umlcompact=4.2cm] at (0,0) {
    \umltitre{\underline{evt-7f3a \,: DetectionEvent}}
    \nodepart{two}
    category = "data\_exfiltration"\\
    severity = HIGH\\
    confidence = 0.88\\
    indicators = \{destIp \,: 198.51.100.77,\\ \phantom{xx} destPort \,: 8443\}
  };
  \node (o2) [umlcompact=3.6cm] at (5,0) {
    \umltitre{\underline{srv-app-01 \,: Asset}}
    \nodepart{two}
    criticality = 4\\
    zone = "dmz"
  };
  \node (o3) [umlcompact=4.4cm] at (10.4,0) {
    \umltitre{\underline{inc-2ed5 \,: Incident}}
    \nodepart{two}
    status = CONTAINED\\
    attackCode = "B4"\\
    dangerousness = 7.4\\
    priorityRank = 2\\
    riskScore() = 71.5
  };
  \node (o4) [umlcompact=4.4cm] at (10.4,-5) {
    \umltitre{\underline{dec-91c0 \,: Decision}}
    \nodepart{two}
    outcome = AUTONOMOUS\_EXECUTION\\
    rationale = "B4 --- exfiltration ;\\ \phantom{xx} 1 action autoris\'ee"
  };
  \node (o5) [umlcompact=4.6cm] at (4.6,-5) {
    \umltitre{\underline{tr-91c0 \,: DecisionTrace}}
    \nodepart{two}
    playbookId = "pb-b4-exfiltration"\\
    groundingScore = 0.81\\
    consideredActions = \{cut\_egress,\\ \phantom{xx} disable\_account\}\\
    rejectedActions = \{disable\_account \,:\\ \phantom{xx} "irr\'eversible --- EF-14"\}
  };
  \node (o6) [umlcompact=4.2cm] at (-0.6,-5) {
    \umltitre{\underline{act-4b12 \,: ActionResult}}
    \nodepart{two}
    status = EXECUTED\\
    rollbackToken = "eg-4b12-77"\\
    durationMs() = 412
  };
  \node (o7) [umlcompact=4.2cm] at (-0.6,-9.6) {
    \umltitre{\underline{sp-cut \,: ActionSpec}}
    \nodepart{two}
    verb = "cut\_egress\_connection"\\
    actuator = "network"\\
    reversibility = PARTIALLY\_REVERSIBLE\\
    rollbackVerb = "restore\_egress"\\
    blastRadius = 1
  };

  \draw[association] (o1) -- (o2);
  \draw[association] (o1) -- (o3);
  \draw[association] (o3) -- (o4);
  \draw[association] (o4) -- (o5);
  \draw[association] (o3.west) to[out=180,in=90] (o6.north);
  \draw[association] (o6) -- (o7);
\end{diagramme}
\caption{Diagramme d'objets --- instantan\'e apr\`es traitement d'une exfiltration}
\label{fig:objets}
\end{figure}

\subsection{Contrat des actionneurs}
Le moteur ne conna\^it aucune technologie d'\'equipement : il ne conna\^it qu'un
contrat. La figure~\ref{fig:actionneurs} le pr\'esente sous forme d'interface et de
r\'ealisations. Deux clauses de ce contrat sont non n\'egociables et expliquent sa
forme.

\begin{description}
  \item[Idempotence.] \texttt{execute()} appel\'ee deux fois avec les m\^emes
    param\`etres produit le m\^eme \'etat final. Sans cette clause, un rejeu de la file
    du mode d\'egrad\'e (ENF-05) doublerait les gestes.
  \item[Remise d'un jeton.] \texttt{execute()} rend un \texttt{rollbackToken}
    d\`es lors que l'action se d\'eclare annulable. C'est ce jeton, et lui seul, qui
    rend la boucle de contr\^ole (EF-25) op\'erante : l'annulation ne se d\'eduit pas de
    l'action, elle s'appuie sur une information produite \`a l'ex\'ecution.
\end{description}

\begin{figure}[H]
\centering
\begin{diagramme}
  \node (iface) [umlclasse=5cm] at (5.5,0) {
    \umltitres{interface}{Actuator}
    \nodepart{two}
    + name \,: String
    \nodepart{three}
    + supports(verb) \,: Boolean\\
    + execute(spec) \,: ActuationOutcome\\
    + rollback(verb, token) \,: ActuationOutcome
  };
  \node (out) [umlcompact=3.8cm] at (12.2,0) {
    \umltitres{value object}{ActuationOutcome}
    \nodepart{two}
    - success \,: Boolean\\
    - rollbackToken \,: String [0..1]\\
    - detail \,: Map
  };
  \node (fw)  [umlsimple=2.4cm] at (-0.4,-4.4) {\umltitre{FirewallActuator}};
  \node (edr) [umlsimple=2.4cm] at (2.4,-4.4) {\umltitre{EdrActuator}};
  \node (iam) [umlsimple=2.4cm] at (5.2,-4.4) {\umltitre{IamActuator}};
  \node (net) [umlsimple=2.4cm] at (8.0,-4.4) {\umltitre{NetworkActuator}};
  \node (dns) [umlsimple=2.4cm] at (10.8,-4.4) {\umltitre{DnsActuator}};
  \node (etc) [umlsimple=2.4cm] at (13.6,-4.4) {\umltitre{\ldots\ (12 au total)}};

  \node (sim) [umlcompact=4.6cm] at (5.5,-8.4) {
    \umltitres{decorator}{SimulationActuator}
    \nodepart{two}
    Enveloppe tout actionneur. En posture \emph{simulation},
    aucun geste n'atteint l'\'equipement r\'eel ; le jeton et
    la trace sont produits \`a l'identique.
  };

  \foreach \n in {fw,edr,iam,net,dns,etc}{ \draw[realisation] (\n.north) -- (iface.south); }
  \draw[dependance] (iface.east) -- node[role]{rend} (out.west);
  \draw[realisation] (sim.north) -- (iface.south);
\end{diagramme}
\caption{Contrat des actionneurs --- patron Strat\'egie et d\'ecorateur de simulation}
\label{fig:actionneurs}
\end{figure}

\subsection{Classes du confinement hors catalogue}
La figure~\ref{fig:classes-repli} pr\'esente les classes engag\'ees lorsque aucun
playbook ne s'applique. Leur particularit\'e tient \`a la source du raisonnement : le
planificateur de repli ne part pas d'un \emph{type d'attaque suppos\'e} --- ce serait
une hypoth\`ese --- mais d'un \emph{indicateur constat\'e} --- ce qui est un fait.
Cette distinction ne l\`eve pas l'exigence de non-invention (EF-04) : elle en change
le fondement.

\begin{figure}[H]
\centering
\begin{diagramme}
  \node (plan) [umlclasse=5.2cm] at (0,0) {
    \umltitre{FallbackPlanner}
    \nodepart{two}
    - catalogue \,: ReversibilityCatalog\\
    - rayonMaximalAutonome \,: Integer = 2\\
    - reseauxInternes \,: IpNetwork [*]
    \nodepart{three}
    + plan(event) \,: FallbackPlan\\
    - candidats(event) \,: Suggestion [*]\\
    - estExterne(adresse) \,: Boolean
  };
  \node (fplan) [umlclasse=4.8cm] at (7.6,0) {
    \umltitre{FallbackPlan}
    \nodepart{two}
    - autonomous \,: Suggestion [*]\\
    - requiresConfirmation \,: Suggestion [*]\\
    - observations \,: String [*]
    \nodepart{three}
    + isEmpty() \,: Boolean\\
    + toDict() \,: Map
  };
  \node (sug) [umlclasse=4.6cm] at (13.8,0) {
    \umltitres{value object}{Suggestion}
    \nodepart{two}
    - spec \,: ActionSpec\\
    - basis \,: String\\
    - autonomous \,: Boolean\\
    - residualEffect \,: String
  };
  \node (cat) [umlcompact=4.6cm] at (0,-6) {
    \umltitre{ReversibilityCatalog}
    \nodepart{two}
    + reversibilityOf(key) \,: Reversibility\\
    + rollbackVerbOf(key) \,: String [0..1]\\
    + typicalBlastRadius(key) \,: Integer
  };
  \node (regle) [note=6.4cm] at (8.4,-6) {
    \textbf{\{r\`egle d'autonomie du repli\}}\\[2pt]
    Une suggestion n'est engag\'ee seule que si\\
    \texttt{reversibility = REVERSIBLE} $\wedge$\\
    \texttt{rollbackVerb <> null} $\wedge$\\
    \texttt{typicalBlastRadius <= 2}.\\[3pt]
    Toute autre suggestion rejoint
    \texttt{requiresConfirmation} et d\'eclenche une alerte
    persistante \`a destination de l'analyste (EF-28).
  };

  \draw[composition] (plan.east) -- node[mult, near start]{1} node[mult, near end]{1} (fplan.west);
  \draw[composition] (fplan.east) -- node[mult, near start]{1} node[mult, near end]{0..*} (sug.west);
  \draw[association] (plan.south) -- node[mult, near end]{1} (cat.north);
  \draw[dependance]  (regle.north west) -- (fplan.south);
\end{diagramme}
\caption{Classes du confinement hors catalogue}
\label{fig:classes-repli}
\end{figure}

%%%==========================================================================
\section{Vue dynamique}
%%%==========================================================================

\subsection{Diagramme de s\'equence syst\`eme}
Le diagramme de s\'equence syst\`eme traite la plateforme comme une bo\^ite noire et ne
retient que les \'echanges avec l'ext\'erieur. Son int\'er\^et ici est de rendre visible,
d'un seul coup d'\oe il, la propri\'et\'e qui distingue le syst\`eme des plateformes
d'orchestration usuelles : \textbf{aucun message n'est adress\'e \`a l'analyste entre
la r\'eception de l'\'ev\'enement et l'action sur l'\'equipement}. Le seul message qu'il
re\c coit est post\'erieur \`a l'ex\'ecution.

\begin{figure}[H]
\centering
\begin{diagramme}
  \seqvie{src}{0}{-7.2}{Source de\\d\'etection}
  \seqvie{sys}{4.4}{-7.2}{\underline{\,:CIRTDEFENSE\,}}
  \seqvie{eqp}{9.0}{-7.2}{\'Equipement\\cible}
  \seqvie{ana}{13.0}{-7.2}{Analyste}

  \seqact{4.4}{-1.0}{-6.4}
  \seqmsg{0}{4.4}{-1.0}{alerte brute}
  \seqauto{4.4}{-1.6}{normaliser, d\'edoublonner, corr\'eler}
  \seqauto{4.4}{-2.6}{ancrer, classer, planifier}
  \seqauto{4.4}{-3.6}{filtrer par r\'eversibilit\'e puis par politique}
  \seqact{9.0}{-4.5}{-5.0}
  \seqmsg{4.4}{9.0}{-4.5}{ex\'ecuter le geste correctif}
  \seqrep{9.0}{4.4}{-5.0}{jeton d'annulation}
  \seqmsg{4.4}{13.0}{-5.7}{notification a posteriori}
  \seqauto{4.4}{-6.3}{inscrire au journal}

  \node[note=3.8cm, anchor=north west] at (0.1,-7.7)
    {Le d\'elai entre le premier et le dernier message ne d\'epend
     d'aucune disponibilit\'e humaine. C'est l\`a toute la diff\'erence
     avec une plateforme qui recommande.};
\end{diagramme}
\caption{Diagramme de s\'equence syst\`eme --- r\'eponse autonome (UC-02)}
\label{fig:sequence-systeme}
\end{figure}

\subsection{Diagramme de s\'equence d\'etaill\'e --- r\'eponse autonome}
La figure~\ref{fig:sequence-detail} ouvre la bo\^ite noire. On y suit la
\emph{cha\^ine de responsabilit\'e} : chaque garde peut interrompre le traitement, et
l'interruption produit toujours une inscription motiv\'ee, jamais un silence.

\begin{sidewaysfigure}
\centering
\ajusterpaysage{%
\begin{diagramme}
  \seqvie{src}{0}{-13.4}{Source}
  \seqvie{ada}{2.5}{-13.4}{\underline{\,:Moteur\,}}
  \seqvie{enr}{5.0}{-13.4}{\underline{\,:Enrichisseur\,}}
  \seqvie{pla}{7.5}{-13.4}{\underline{\,:Planificateur\,}}
  \seqvie{cat}{10.0}{-13.4}{\underline{\,:CatalogueR\'ev.\,}}
  \seqvie{pol}{12.5}{-13.4}{\underline{\,:Politique\,}}
  \seqvie{cbk}{15.0}{-13.4}{\underline{\,:CoupeCircuit\,}}
  \seqvie{exe}{17.5}{-13.4}{\underline{\,:Ex\'ecuteur\,}}
  \seqvie{act}{20.0}{-13.4}{\underline{\,:Actionneur\,}}
  \seqvie{jou}{22.5}{-13.4}{\underline{\,:Journal\,}}

  \seqact{2.5}{-1.0}{-12.9}
  \seqmsg{0}{2.5}{-1.0}{charge utile brute}
  \seqauto{2.5}{-1.5}{normaliser via l'adaptateur $\rightarrow$ DetectionEvent (EF-18)}
  \seqauto{2.5}{-2.3}{empreinte, d\'edoublonnage (EF-19)}
  \seqauto{2.5}{-3.1}{corr\'eler $\rightarrow$ Incident (EF-20)}
  \seqmsg{2.5}{22.5}{-3.7}{inscrire \guillemotleft event.ingested\guillemotright}

  \seqact{5.0}{-4.2}{-5.2}
  \seqmsg{2.5}{5.0}{-4.2}{enrichir(incident)}
  \seqauto{5.0}{-4.7}{rechercher au fonds documentaire (EF-03)}
  \seqrep{5.0}{2.5}{-5.2}{Contexte \{ancr\'e, score\}}

  \seqfrag{1.3}{-5.7}{23.4}{-7.2}{alt}{contexte non ancr\'e --- EF-04}
  \node[font=\tiny, anchor=west] at (3.4,-6.7)
    {refus d'agir motiv\'e $\rightarrow$ plan de repli (fig.~\ref{fig:sequence-repli}) $\rightarrow$ alerte d'abstention. \textbf{Aucune action ex\'ecut\'ee.}};

  \seqact{7.5}{-7.5}{-9.2}
  \seqmsg{2.5}{7.5}{-7.5}{planifier(incident, contexte)}
  \seqauto{7.5}{-8.0}{classer au catalogue m\'etier}
  \seqmsg{7.5}{10.0}{-8.7}{r\'eversibilit\'e ?}
  \seqact{10.0}{-8.7}{-9.2}
  \seqrep{10.0}{7.5}{-9.2}{\{r\'eversible, verbe, rayon\}}

  \seqact{12.5}{-9.7}{-10.2}
  \seqmsg{7.5}{12.5}{-9.7}{\'evaluer(actions) (EF-15)}
  \seqrep{12.5}{7.5}{-10.2}{PolicyVerdict [*]}

  \seqact{15.0}{-10.6}{-11.0}
  \seqmsg{7.5}{15.0}{-10.6}{autonomie autoris\'ee ? (EF-26)}
  \seqrep{15.0}{7.5}{-11.0}{ferm\'e}

  \seqact{17.5}{-11.4}{-12.6}
  \seqmsg{7.5}{17.5}{-11.4}{ex\'ecuter(actions autoris\'ees) (EF-07)}
  \seqact{20.0}{-11.9}{-12.3}
  \seqmsg{17.5}{20.0}{-11.9}{execute(spec)}
  \seqrep{20.0}{17.5}{-12.3}{outcome \{jeton\}}
  \seqmsg{17.5}{22.5}{-12.9}{inscrire \guillemotleft action.executed\guillemotright}
\end{diagramme}}
\caption{Diagramme de s\'equence d\'etaill\'e --- cha\^ine de d\'ecision et d'ex\'ecution autonome}
\label{fig:sequence-detail}
\end{sidewaysfigure}

\subsection{Diagramme de s\'equence --- confinement hors catalogue}
Ce sc\'enario est celui qui distingue le plus nettement le syst\`eme d'une plateforme
\`a playbooks. Lorsque le contexte n'est pas document\'e, la r\'eponse usuelle est le
silence. Ici, le moteur bascule vers le planificateur de repli, qui raisonne sur
les indicateurs et non sur une hypoth\`ese de type d'attaque. Le fragment
\texttt{loop} de la figure~\ref{fig:sequence-repli} porte le c\oe ur du m\'ecanisme :
chaque indicateur exploitable produit au plus un geste, et ce geste est aussit\^ot
class\'e selon qu'il est annulable ou non.

\begin{sidewaysfigure}
\centering
\ajusterpaysage{%
\begin{diagramme}
  \seqvie{mot}{0}{-13.0}{\underline{\,:Moteur\,}}
  \seqvie{rep}{3.6}{-13.0}{\underline{\,:PlanificateurRepli\,}}
  \seqvie{cat}{7.6}{-13.0}{\underline{\,:CatalogueR\'ev.\,}}
  \seqvie{pol}{11.2}{-13.0}{\underline{\,:Politique\,}}
  \seqvie{exe}{14.6}{-13.0}{\underline{\,:Ex\'ecuteur\,}}
  \seqvie{ale}{18.0}{-13.0}{\underline{\,:Alerte\,}}
  \seqvie{jou}{21.4}{-13.0}{\underline{\,:Journal\,}}

  \seqact{0}{-0.9}{-12.6}
  \seqauto{0}{-1.1}{contexte non ancr\'e (EF-04)}

  \seqact{3.6}{-2.0}{-8.4}
  \seqmsg{0}{3.6}{-2.0}{plan(event)}
  \seqauto{3.6}{-2.5}{extraire les indicateurs observ\'es}

  \seqfrag{2.9}{-3.2}{9.4}{-7.3}{loop}{pour chaque indicateur exploitable}
  \seqauto{3.6}{-4.0}{d\'eduire un geste de confinement}
  \seqmsg{3.6}{7.6}{-4.9}{m\'etadonn\'ees(geste)}
  \seqact{7.6}{-4.9}{-5.4}
  \seqrep{7.6}{3.6}{-5.4}{\{r\'eversibilit\'e, verbe, rayon\}}
  \seqauto{3.6}{-6.1}{r\'eversible $\wedge$ rayon $\leq$ 2 ?}
  \node[font=\tiny, anchor=west, align=left] at (4.6,-6.9)
    {oui $\rightarrow$ \texttt{autonomous} \quad non $\rightarrow$ \texttt{requiresConfirmation}};

  \seqrep{3.6}{0}{-8.4}{FallbackPlan}

  \seqact{11.2}{-9.0}{-9.5}
  \seqmsg{0}{11.2}{-9.0}{\'evaluer(gestes autonomes)}
  \seqrep{11.2}{0}{-9.5}{verdicts --- \textbf{la politique s'applique aussi au repli}}

  \seqact{14.6}{-10.1}{-10.6}
  \seqmsg{0}{14.6}{-10.1}{ex\'ecuter(gestes autoris\'es)}
  \seqrep{14.6}{0}{-10.6}{jetons d'annulation}

  \seqact{18.0}{-11.3}{-11.8}
  \seqmsg{0}{18.0}{-11.3}{alerte persistante(gestes \`a effet durable) (EF-28)}
  \seqmsg{0}{21.4}{-12.4}{inscrire \guillemotleft decision.made \{fallback\}\guillemotright}
\end{diagramme}}
\caption{Diagramme de s\'equence --- confinement d'une menace hors catalogue}
\label{fig:sequence-repli}
\end{sidewaysfigure}

\subsection{Diagramme de s\'equence --- boucle de contr\^ole post-action}
La boucle de contr\^ole traite le mode de d\'efaillance le plus subtil : le syst\`eme a
raison sur la menace, mais son geste nuit \`a la cible. Un blocage l\'egitime peut
couper un flux m\'etier ; une isolation justifi\'ee peut rendre un service
indisponible. Le principe est celui d'une mesure diff\'erentielle : une r\'ef\'erence
est prise \emph{avant} l'action, une mesure \emph{apr\`es}, et l'\'ecart d\'ecide.

\begin{figure}[H]
\centering
\ajusterportrait{%
\begin{diagramme}
  \seqvie{ord}{0}{-10.6}{Ordonnanceur}
  \seqvie{bcl}{3.4}{-10.6}{\underline{\,:BouclePostAction\,}}
  \seqvie{son}{7.2}{-10.6}{\underline{\,:Sonde\,}}
  \seqvie{rbk}{10.6}{-10.6}{\underline{\,:Rollback\,}}
  \seqvie{act}{14.0}{-10.6}{\underline{\,:Actionneur\,}}
  \seqvie{jou}{17.4}{-10.6}{\underline{\,:Journal\,}}

  \seqact{3.4}{-1.0}{-10.2}
  \seqmsg{0}{3.4}{-1.0}{d\'eclencher le cycle}

  \seqfrag{2.7}{-1.7}{18.4}{-9.6}{loop}{pour chaque action ex\'ecut\'ee et annulable}

  \seqact{7.2}{-2.6}{-3.1}
  \seqmsg{3.4}{7.2}{-2.6}{mesurer(cible)}
  \seqrep{7.2}{3.4}{-3.1}{indicateurs de service}

  \seqfrag{3.0}{-3.6}{18.1}{-5.0}{alt}{aucune r\'ef\'erence ant\'erieure}
  \node[font=\tiny, anchor=west] at (3.3,-4.5)
    {abstention motiv\'ee et inscrite --- sans mesure ant\'erieure, la comparaison n'aurait aucun sens};

  \seqauto{3.4}{-5.6}{comparer \`a la r\'ef\'erence prise avant l'action}

  \seqfrag{3.0}{-6.4}{18.1}{-9.3}{opt}{d\'egradation $>$ seuil}
  \seqact{10.6}{-7.2}{-8.8}
  \seqmsg{3.4}{10.6}{-7.2}{annuler(action, jeton)}
  \seqact{14.0}{-7.7}{-8.2}
  \seqmsg{10.6}{14.0}{-7.7}{rollback(verbe, jeton)}
  \seqrep{14.0}{10.6}{-8.2}{r\'esultat}
  \seqmsg{10.6}{17.4}{-8.9}{inscrire \guillemotleft rollback.completed\guillemotright}

  \node[note=6.2cm, anchor=north west] at (0,-11.1)
    {\textbf{Limite connue et assum\'ee.} La r\'ef\'erence prise avant l'action est
     conserv\'ee en m\'emoire vive. Apr\`es un red\'emarrage, la boucle s'abstient sur
     les actions ant\'erieures : elle pr\'ef\`ere ne rien faire plut\^ot que comparer
     \`a une r\'ef\'erence qu'elle n'a pas prise. Le report de cette r\'ef\'erence en base
     figure aux perspectives de la conclusion g\'en\'erale.};
\end{diagramme}}
\caption{Diagramme de s\'equence --- boucle de contr\^ole post-action (EF-25)}
\label{fig:sequence-boucle}
\end{figure}

\subsection{Diagramme d'activit\'e}
Le diagramme d'activit\'e de la figure~\ref{fig:activite} donne la vue algorithmique
du traitement complet, couloirs compris. Sa lecture appelle deux remarques.

D'abord, \textbf{tout chemin de refus aboutit \`a une inscription et \`a une alerte},
jamais \`a une sortie silencieuse. C'est une exigence de conception : un syst\`eme
autonome qui renonce sans le dire est plus dangereux qu'un syst\`eme qui n'a rien
tent\'e, car l'exploitant croit la menace trait\'ee.

Ensuite, aucun flot de contr\^ole n'atteint un acteur humain avant l'ex\'ecution. Les
seuls messages qui lui parviennent partent des activit\'es situ\'ees \emph{apr\`es}
l'action sur l'\'equipement.

\begin{figure}[p]
\centering
\ajusterportrait{%
\begin{diagramme}
  %%% ---------------------------------------------- couloir 1 : acquisition
  \node (deb)  [etatinit] at (0,0) {};
  \node (a1)   [action=2.7cm] at (0,-1.5) {Normaliser l'\'ev\'enement (EF-18)};
  \node (d1)   [decision] at (0,-3.3) {d\'ej\`a vu ?};
  \node (rien) [etatfin] at (-2.6,-3.3) {};
  \node (a2)   [action=2.7cm] at (0,-5.1) {Corr\'eler en incident (EF-20)};

  %%% -------------------------------------------------- couloir 2 : analyse
  \node (a3)   [action=2.7cm] at (4.4,-5.1) {Enrichir au fonds documentaire (EF-03)};
  \node (d2)   [decision] at (4.4,-6.9) {contexte ancr\'e ?};
  \node (rep)  [action=2.9cm, draw=umlaccent] at (8.2,-6.9)
        {Confinement de repli (EF-04)\\ \tiny fig.~\ref{fig:activite-repli}};
  \node (a4)   [action=2.7cm] at (4.4,-8.7) {Classer au catalogue m\'etier};
  \node (a5)   [action=2.7cm] at (4.4,-10.5) {Planifier les actions (EF-05, EF-06)};

  %%% --------------------------------------------------- couloir 3 : gardes
  \node (a6)   [action=2.9cm] at (11.6,-10.5) {Filtrer par r\'eversibilit\'e (EF-14)};
  \node (d3)   [decision] at (11.6,-12.1) {une action au moins ?};
  \node (a7)   [action=2.9cm] at (11.6,-13.7) {\'Evaluer la politique (EF-15)};
  \node (d4)   [decision] at (11.6,-15.3) {autoris\'ee ?};
  \node (d5)   [decision] at (11.6,-16.9) {coupe-circuit ferm\'e ?};
  \node (x1)   [action=2.9cm, draw=umlalerte] at (11.6,-19.0)
        {Inscrire un refus motiv\'e et alerter : menace non contenue};
  \node (finx) [etatfin] at (11.6,-20.7) {};

  %%% ------------------------------------ couloir 4 : execution et retour
  \node (e1)   [action=2.9cm] at (16.0,-16.9) {Ex\'ecuter aupr\`es des actionneurs (EF-07)};
  \node (e2)   [action=2.9cm] at (16.0,-18.5) {Recueillir les jetons d'annulation};
  \node (e3)   [action=2.9cm] at (16.0,-20.1) {Armer la boucle de contr\^ole (EF-25)};
  \node (e4)   [action=2.9cm] at (16.0,-21.7) {Inscrire au journal};
  \node (e5)   [action=2.9cm] at (16.0,-23.3) {Notifier l'analyste (EF-13)};
  \node (fin)  [etatfin] at (16.0,-24.8) {};

  %%% ------------------------------------------------------------- flots
  \draw[flot]   (deb) -- (a1);
  \draw[flot]   (a1) -- (d1);
  \draw[flotko] (d1) -- node[garde, above]{oui} node[garde, below]{\'ecart\'e (EF-19)} (rien);
  \draw[flot]   (d1) -- node[garde, right]{non} (a2);
  \draw[flot]   (a2) -- (a3);
  \draw[flot]   (a3) -- (d2);
  \draw[flotok] (d2) -- node[garde, right]{oui} (a4);
  \draw[flotko] (d2) -- node[garde, above]{non} (rep);
  \draw[flot]   (a4) -- (a5);
  \draw[flot]   (a5) -- (a6);
  \draw[flot]   (a6) -- (d3);
  \draw[flotok] (d3) -- node[garde, right]{oui} (a7);
  \draw[flot]   (a7) -- (d4);
  \draw[flotok] (d4) -- node[garde, right]{oui} (d5);
  \draw[flotok] (d5) -- node[garde, above]{ferm\'e} (e1);
  \draw[flotko] (d3.west) -- ++(-1.0,0) |- node[garde, pos=0.15, left]{non} (x1.west);
  \draw[flotko] (d4.west) -- ++(-0.5,0) |- node[garde, pos=0.15, left]{non} ([yshift=3mm]x1.west);
  \draw[flotko] (d5) -- node[garde, right]{ouvert} (x1);
  \draw[flot]   (x1) -- (finx);
  \draw[flot]   (e1) -- (e2);
  \draw[flot]   (e2) -- (e3);
  \draw[flot]   (e3) -- (e4);
  \draw[flot]   (e4) -- (e5);
  \draw[flot]   (e5) -- (fin);
  \draw[flot]   (rep.east) -- (13.6,-6.9) -- (13.6,-18.5) -- (e2.west);

  %%% ---------------------------------------------------------- couloirs
  \begin{scope}[on background layer]
    \node (c1) [couloir, fit=(deb)(a1)(d1)(rien)(a2)] {};
    \node (c2) [couloir, fit=(a3)(d2)(rep)(a4)(a5)] {};
    \node (c3) [couloir, fit=(a6)(d3)(a7)(d4)(d5)(x1)(finx)] {};
    \node (c4) [couloir, fit=(e1)(e2)(e3)(e4)(e5)(fin)] {};
  \end{scope}
  \node[etiquettecouloir, anchor=south] at (c1.west) {Acquisition};
  \node[etiquettecouloir, anchor=south] at (c2.west) {Analyse};
  \node[etiquettecouloir, anchor=south] at (c3.west) {Gardes};
  \node[etiquettecouloir, anchor=south] at (c4.west) {Ex\'ecution et retour};
\end{diagramme}}
\caption{Diagramme d'activit\'e avec couloirs --- cha\^ine principale de traitement}
\label{fig:activite}
\end{figure}

La branche de repli, r\'esum\'ee \`a une activit\'e unique sur la
figure~\ref{fig:activite}, est d\'etaill\'ee \`a la figure~\ref{fig:activite-repli}. Le
n\oe ud de d\'ecision qu'elle contient est le point exact o\`u se joue l'arbitrage
demand\'e par le CIRT : ce qui est annulable part sans attendre, ce qui laisse une
trace durable attend une d\'ecision humaine.

\begin{figure}[H]
\centering
\ajusterportrait{%
\begin{diagramme}
  \node (i)   [etatinit] at (0,0) {};
  \node (b1)  [action=3.2cm] at (0,-1.6) {Extraire les indicateurs observ\'es de l'\'ev\'enement};
  \node (d0)  [decision] at (0,-3.6) {au moins un indicateur exploitable ?};
  \node (z1)  [action=3.2cm, draw=umlalerte] at (5.4,-3.6) {Alerte d'abstention motiv\'ee};
  \node (fz)  [etatfin] at (9.0,-3.6) {};
  \node (b2)  [action=3.2cm] at (0,-5.8) {D\'eduire un geste de confinement par indicateur};
  \node (b3)  [action=3.2cm] at (0,-7.8) {Consulter le catalogue de r\'eversibilit\'e};
  \node (d6)  [decision] at (0,-10.0) {r\'eversible et rayon $\leq$ 2 ?};
  \node (g1)  [action=3.2cm] at (-3.4,-12.2) {Soumettre \`a la politique (EF-15)};
  \node (g2)  [action=3.2cm] at (4.0,-12.2) {Marquer \guillemotleft en attente de d\'ecision humaine\guillemotright};
  \node (h1)  [action=3.2cm] at (-3.4,-14.2) {Ex\'ecuter sans attendre};
  \node (h2)  [action=3.2cm] at (4.0,-14.2) {\'Emettre une alerte persistante (EF-28)};
  \node (j)   [barre=8.4cm] at (0.3,-16.0) {};
  \node (k1)  [action=3.4cm] at (0.3,-17.6) {Inscrire au journal le partage entre gestes engag\'es et gestes en attente};
  \node (k2)  [action=3.4cm] at (0.3,-19.6) {Ouvrir une demande de qualification (EF-29)};
  \node (fin) [etatfin] at (0.3,-21.2) {};

  \draw[flot]   (i) -- (b1);
  \draw[flot]   (b1) -- (d0);
  \draw[flotko] (d0) -- node[garde, above]{non} (z1);
  \draw[flot]   (z1) -- (fz);
  \draw[flotok] (d0) -- node[garde, right]{oui} (b2);
  \draw[flot]   (b2) -- (b3);
  \draw[flot]   (b3) -- (d6);
  \draw[flotok] (d6.west) to[out=180,in=90] node[garde, left]{oui} (g1.north);
  \draw[flotko] (d6.east) to[out=0,in=90] node[garde, right]{non} (g2.north);
  \draw[flot]   (g1) -- (h1);
  \draw[flot]   (g2) -- (h2);
  \draw[flot]   (h1.south) -- (j.north west);
  \draw[flot]   (h2.south) -- (j.north east);
  \draw[flot]   (j) -- (k1);
  \draw[flot]   (k1) -- (k2);
  \draw[flot]   (k2) -- (fin);

  \node[note=4.6cm, anchor=north west] at (5.6,-6.0)
    {Le geste n'est jamais d\'eduit d'un \emph{type d'attaque suppos\'e} mais
     d'un \emph{indicateur constat\'e}. L'exigence de non-invention (EF-04) n'est
     donc pas lev\'ee : son fondement change.};
\end{diagramme}}
\caption{Diagramme d'activit\'e --- branche de confinement hors catalogue}
\label{fig:activite-repli}
\end{figure}

\subsection{Diagrammes d'\'etats-transitions}
\label{sec:etats}

\subsubsection{Cycle de vie d'un incident}
La figure~\ref{fig:etats-incident} pr\'esente la machine \`a \'etats de l'incident. Ce
qui la caract\'erise se lit \emph{en creux} : \textbf{il n'existe aucun \'etat
\og en attente de validation \fg{}}. C'est la traduction structurelle du choix
d'autonomie --- un tel \'etat aurait r\'eintroduit, dans le mod\`ele lui-m\^eme, la
d\'ependance \`a la disponibilit\'e humaine que le projet entend lever.

\begin{figure}[H]
\centering
\ajusterportrait{%
\begin{diagramme}
  \node (i)   [etatinit] at (0,0) {};
  \node (op)  [etat] at (4.6,0) {ouvert};
  \node (co)  [etat] at (9.6,0) {confin\'e};
  \node (rb)  [etat] at (9.6,-3.4) {annul\'e};
  \node (cl)  [etat] at (14.4,0) {clos};
  \node (f)   [etatfin] at (16.8,0) {};

  \draw[transition] (i) -- node[garde, above]{\'ev\'enement corr\'el\'e} (op);
  \draw[transition] (op) -- node[garde, above, align=center]{actions ex\'ecut\'ees\\ \tiny[au moins une r\'eussie]} (co);
  \draw[transition] (co) -- node[garde, right, align=left]{d\'egradation d\'etect\'ee\\ \tiny/ annuler les actions} (rb);
  \draw[transition] (co) -- node[garde, above]{d\'elai de veille \'ecoul\'e} (cl);
  \draw[transition] (rb.east) to[out=0,in=270] node[garde, below right]{constat de l'analyste} (cl.south);
  \draw[transition] (cl) -- (f);

  \path[transition] (op) edge[loop above] node[garde, align=center]
        {\'ev\'enement corr\'el\'e suppl\'ementaire\\ \tiny/ absorber} ();
  \path[transition] (op) edge[loop below] node[garde, align=center]
        {aucune action possible\\ \tiny/ inscrire un refus motiv\'e} ();

  \node[note=5.0cm, anchor=north west] at (11.8,-2.4)
    {Aucun \'etat \og en attente de validation \fg{} ne figure sur ce diagramme.
     Son absence est un r\'esultat de conception, non un oubli : c'est elle qui
     garantit que le d\'elai de r\'eponse ne d\'epend d'aucune disponibilit\'e humaine.};
\end{diagramme}}
\caption{Diagramme d'\'etats-transitions --- cycle de vie d'un incident}
\label{fig:etats-incident}
\end{figure}

\subsubsection{Cycle de vie d'une action}
La machine \`a \'etats de l'action, figure~\ref{fig:etats-action}, porte les huit
\'etats de l'\'enum\'eration \texttt{ActionStatus}. Trois d'entre eux sont des refus
ant\'erieurs \`a toute ex\'ecution --- \emph{bloqu\'ee par la politique}, \emph{bloqu\'ee par
le coupe-circuit} --- ou un \'echec de rattrapage --- \emph{\'echec d'annulation}. Ce
dernier est le seul \'etat du syst\`eme qui exige une intervention humaine imm\'ediate :
une action a \'et\'e men\'ee et ne peut plus \^etre d\'efaite.

\begin{figure}[H]
\centering
\begin{diagramme}[node distance=8mm]
  \node (i)   [etatinit] at (0,0) {};
  \node (pl)  [etat] at (2.4,0) {planifi\'ee};
  \node (bp)  [etat, draw=umlalerte] at (2.4,-2.4) {bloqu\'ee\\\tiny par la politique};
  \node (bc)  [etat, draw=umlalerte] at (2.4,-4.4) {bloqu\'ee\\\tiny par le coupe-circuit};
  \node (ex)  [etat] at (6.4,0) {en cours};
  \node (ok)  [etat] at (10.4,0) {ex\'ecut\'ee};
  \node (ko)  [etat, draw=umlalerte] at (6.4,-2.4) {en \'echec};
  \node (an)  [etat] at (10.4,-2.4) {annul\'ee};
  \node (ae)  [etat, draw=umlalerte, line width=1pt] at (10.4,-4.6) {\'echec\\d'annulation};
  \node (f)   [etatfin] at (13.6,-1.2) {};

  \draw[transition] (i) -- (pl);
  \draw[transition] (pl) -- node[garde, above]{autoris\'ee} (ex);
  \draw[transition] (pl) -- node[garde, right]{refus\'ee (EF-15)} (bp);
  \draw[transition] (pl.west) to[out=180,in=180] node[garde, left]{autonomie suspendue (EF-26)} (bc.west);
  \draw[transition] (ex) -- node[garde, above]{succ\`es} (ok);
  \draw[transition] (ex) -- node[garde, right]{\'echec de l'actionneur} (ko);
  \draw[transition] (ok) -- node[garde, left, align=right]{d\'egradation constat\'ee\\ \tiny[jeton disponible] / rollback} (an);
  \draw[transition] (an) -- node[garde, right]{le rollback \'echoue} (ae);
  \draw[transition] (ok.east) -- ++(1.4,0) |- (f.north);
  \draw[transition] (an.east) -- ++(1.4,0) -- (f.west);
  \draw[transition] (ko.south) |- ([yshift=-1.2cm]an.south) -| (f.south);
\end{diagramme}
\caption{Diagramme d'\'etats-transitions --- cycle de vie d'une action corrective}
\label{fig:etats-action}
\end{figure}

\subsection{Diagramme de communication}
Le diagramme de communication de la figure~\ref{fig:communication} pr\'esente les
m\^emes \'echanges que la figure~\ref{fig:sequence-detail}, mais organis\'es par
\emph{voisinage} plut\^ot que par chronologie. Il fait ressortir un point que la
lecture s\'equentielle masque : le \emph{moteur} est le seul objet en relation avec
tous les autres. Cette centralit\'e est voulue --- elle concentre en un point unique
la responsabilit\'e d'appliquer les gardes dans le bon ordre --- mais elle d\'esigne
aussi le module dont toute modification demande le plus de prudence.

\begin{figure}[H]
\centering
\begin{diagramme}
  \node (mot) [umlsimple=2.6cm, umlaccentue] at (0,0) {\umltitre{\underline{\,:Moteur\,}}};
  \node (ada) [umlsimple=2.4cm] at (-4.6,2.6) {\umltitre{\underline{\,:Adaptateur\,}}};
  \node (enr) [umlsimple=2.4cm] at (0,3.4) {\umltitre{\underline{\,:Enrichisseur\,}}};
  \node (pla) [umlsimple=2.4cm] at (4.6,2.6) {\umltitre{\underline{\,:Planificateur\,}}};
  \node (pol) [umlsimple=2.4cm] at (5.6,0) {\umltitre{\underline{\,:Politique\,}}};
  \node (cbk) [umlsimple=2.4cm] at (4.6,-2.6) {\umltitre{\underline{\,:CoupeCircuit\,}}};
  \node (exe) [umlsimple=2.4cm] at (0,-3.4) {\umltitre{\underline{\,:Ex\'ecuteur\,}}};
  \node (jou) [umlsimple=2.4cm] at (-4.6,-2.6) {\umltitre{\underline{\,:Journal\,}}};
  \node (rep) [umlsimple=2.4cm] at (-5.6,0) {\umltitre{\underline{\,:Repli\,}}};

  \draw[association] (ada) -- node[role, pos=0.55]{1\,: incident} (mot);
  \draw[association] (mot) -- node[role, pos=0.5]{2\,: enrichir} (enr);
  \draw[association] (mot) -- node[role, pos=0.5]{3\,: planifier} (pla);
  \draw[association] (mot) -- node[role, pos=0.5]{4\,: \'evaluer} (pol);
  \draw[association] (mot) -- node[role, pos=0.5]{5\,: autoris\'e ?} (cbk);
  \draw[association] (mot) -- node[role, pos=0.5]{6\,: ex\'ecuter} (exe);
  \draw[association] (mot) -- node[role, pos=0.5]{7\,: inscrire} (jou);
  \draw[association] (mot) -- node[role, pos=0.5]{2b\,: plan de repli} (rep);
\end{diagramme}
\caption{Diagramme de communication --- voisinage du moteur d'orchestration}
\label{fig:communication}
\end{figure}

%%%==========================================================================
\section{Vue architecturale}
%%%==========================================================================

\subsection{Diagramme de composants}
La figure~\ref{fig:composants} pr\'esente le d\'ecoupage en composants et les
interfaces qu'ils exposent ou requi\`erent. Le \emph{Moteur d'orchestration} y
occupe une position particuli\`ere : il ne fournit qu'une interface --- traiter un
\'ev\'enement --- mais en requiert six. C'est le sympt\^ome, assum\'e, d'un composant qui
coordonne sans rien impl\'ementer lui-m\^eme.

\begin{sidewaysfigure}
\centering
\ajusterpaysage{%
\begin{diagramme}
  %%% ------------------------------------------------------- entrees
  \node (src)  [composant=3.0cm] at (0,2.4)   {\umlstereo{device}\\\textbf{Sources externes}\\\tiny SIEM, sonde r\'eseau, EDR, journaux};
  \node (ueba) [composant=3.0cm] at (0,0)     {\umlstereo{component}\\\textbf{UEBA}\\\tiny profil et \'ecart comportemental};
  \node (infra)[composant=3.0cm] at (0,-2.4)  {\umlstereo{component}\\\textbf{Sondes d'infrastructure}\\\tiny disponibilit\'e, seuils de service};

  %%% ------------------------------------------------------- chaine
  \node (ada)  [composant=3.0cm] at (4.6,0)   {\umlstereo{component}\\\textbf{Adaptateur d'ingestion}\\\tiny normalise, d\'edoublonne, corr\`ele};
  \node (mot)  [composant=3.4cm, umlaccentue] at (9.4,0) {\umlstereo{component}\\\textbf{Moteur d'orchestration}\\\tiny applique les gardes dans l'ordre};
  \node (enr)  [composant=2.9cm] at (9.4,3.4) {\umlstereo{component}\\\textbf{Enrichissement}\\\tiny fonds documentaire, ancrage};
  \node (pol)  [composant=2.9cm] at (13.0,3.4){\umlstereo{component}\\\textbf{Politique compil\'ee}};
  \node (cat)  [composant=2.9cm] at (5.8,3.4) {\umlstereo{component}\\\textbf{Catalogue de r\'eversibilit\'e}};
  \node (cbk)  [composant=2.9cm] at (9.4,-3.4){\umlstereo{component}\\\textbf{Coupe-circuit}};
  \node (rpl)  [composant=2.9cm] at (5.8,-3.4){\umlstereo{component}\\\textbf{Planificateur de repli}};

  %%% ------------------------------------------------------- sorties
  \node (exe)  [composant=2.9cm] at (14.0,0)  {\umlstereo{component}\\\textbf{Ex\'ecuteur}};
  \node (act)  [composant=3.0cm] at (18.4,1.6){\umlstereo{device}\\\textbf{\'Equipements}\\\tiny pare-feu, EDR, annuaire, DNS};
  \node (bcl)  [composant=2.9cm] at (14.0,-3.4){\umlstereo{component}\\\textbf{Boucle post-action}};
  \node (jou)  [composant=2.9cm] at (18.4,-1.6){\umlstereo{component}\\\textbf{Journal cha\^in\'e}};
  \node (not)  [composant=2.9cm] at (18.4,-4.4){\umlstereo{component}\\\textbf{Notification}};
  \node (api)  [composant=2.9cm] at (14.0,-6.4){\umlstereo{component}\\\textbf{API et interface}};

  %%% ------------------------------------------------------- liaisons
  \foreach \n in {src,ueba,infra} { \draw[dependance] (\n.east) -- (ada.west); }
  \draw[dependance] (ada) -- node[role]{\'ev\'enement normalis\'e} (mot);
  \draw[dependance] (mot) -- node[role]{contexte} (enr);
  \draw[dependance] (mot) -- node[role]{verdicts} (pol);
  \draw[dependance] (mot) -- node[role]{r\'eversibilit\'e} (cat);
  \draw[dependance] (mot) -- node[role]{autorisation} (cbk);
  \draw[dependance] (mot) -- node[role]{plan de repli} (rpl);
  \draw[dependance] (mot) -- node[role]{ordres} (exe);
  \draw[dependance] (exe) -- node[role]{gestes} (act);
  \draw[dependance] (exe) -- (jou);
  \draw[dependance] (exe) -- (bcl);
  \draw[dependance] (bcl.east) -- node[role]{annulation} (jou.south west);
  \draw[dependance] (mot.south east) -- (not.north west);
  \draw[dependance] (api.north) -- (exe.south);
  \draw[dependance] (api.east) -- (jou.south);
  \draw[dependance] (infra.south) to[out=270,in=180] node[role, below]{mesure de r\'ef\'erence} (bcl.west);
\end{diagramme}}
\caption{Diagramme de composants}
\label{fig:composants}
\end{sidewaysfigure}

\subsection{Diagramme de d\'eploiement}
La contrainte de souverainet\'e (ENF-01) et celle de sobri\'et\'e d'infrastructure
(ENF-07) r\'eduisent consid\'erablement l'espace des topologies acceptables : ni
service manag\'e, ni d\'ependance \`a un mod\`ele h\'eberg\'e \`a l'ext\'erieur, ni grappe de
serveurs. La figure~\ref{fig:deploiement} pr\'esente la topologie retenue --- un
n\oe ud unique dans le p\'erim\`etre du CIRT, dialoguant avec les \'equipements du parc
supervis\'e sur des protocoles d'administration.

\begin{figure}[H]
\centering
\ajusterportrait{%
\begin{diagramme}
  %%% ---- noeud serveur
  \node (srv) [noeudmat, minimum width=8.6cm, minimum height=6.4cm] at (0,0) {};
  \node[font=\scriptsize\bfseries, anchor=north] at (0,3.0)
       {\umlstereo{device} Serveur CIRT --- p\'erim\`etre souverain};
  \node (app) [artefact, text width=3.4cm] at (-2.1,1.5)
       {\umlstereo{artifact}\\ \texttt{cirtdefense}\\ \tiny service applicatif (ASGI)};
  \node (db)  [artefact, text width=3.4cm] at (-2.1,-0.6)
       {\umlstereo{artifact}\\ \texttt{cirtdefense.db}\\ \tiny base locale, journal cha\^in\'e};
  \node (pbk) [artefact, text width=3.4cm] at (-2.1,-2.4)
       {\umlstereo{artifact}\\ \texttt{playbooks/*.yaml}\\ \tiny proc\'edures de r\'eponse};
  \node (web) [artefact, text width=3.2cm] at (2.2,1.5)
       {\umlstereo{artifact}\\ interface d'exploitation};
  \node (fnd) [artefact, text width=3.2cm] at (2.2,-0.6)
       {\umlstereo{artifact}\\ fonds documentaire local};
  \node (mdl) [artefact, text width=3.2cm] at (2.2,-2.4)
       {\umlstereo{artifact}\\ mod\`ele de langage local\\ \tiny \emph{optionnel --- reformulation seule}};

  %%% ---- postes et equipements
  \node (pos) [noeudmat, text width=2.8cm] at (-6.6,1.2)
       {\umlstereo{device}\\ Poste analyste\\ \tiny navigateur};
  \node (eq1) [noeudmat, text width=2.6cm] at (6.8,2.2) {\umlstereo{device}\\ Pare-feu};
  \node (eq2) [noeudmat, text width=2.6cm] at (6.8,0.6) {\umlstereo{device}\\ Console EDR};
  \node (eq3) [noeudmat, text width=2.6cm] at (6.8,-1.0) {\umlstereo{device}\\ Annuaire};
  \node (eq4) [noeudmat, text width=2.6cm] at (6.8,-2.6) {\umlstereo{device}\\ R\'esolveur DNS};
  \node (col) [noeudmat, text width=2.8cm] at (-6.6,-1.8)
       {\umlstereo{device}\\ Collecteurs\\ \tiny SIEM, sondes};

  \draw[association] (pos.east) -- node[role, above]{HTTPS} (srv.west |- pos);
  \draw[association] (col.east) -- node[role, above]{HTTPS --- flux d'alertes} (srv.west |- col);
  \draw[association] (srv.east |- eq1) -- (eq1.west);
  \draw[association] (srv.east |- eq2) -- (eq2.west);
  \draw[association] (srv.east |- eq3) -- node[role, above]{API d'administration} (eq3.west);
  \draw[association] (srv.east |- eq4) -- (eq4.west);

  \node[note=5.4cm, anchor=north west] at (-8.0,-3.6)
    {\textbf{Aucun lien sortant obligatoire.} Priv\'e de connectivit\'e, le n\oe ud
     conserve la totalit\'e de sa cha\^ine de d\'ecision. Seules les remont\'ees de
     collecteurs distants sont mises en file et rejou\'ees au retour du lien
     (ENF-05).};
\end{diagramme}}
\caption{Diagramme de d\'eploiement --- topologie \`a n\oe ud unique}
\label{fig:deploiement}
\end{figure}

%%%==========================================================================
\section{Conception de la persistance}
%%%==========================================================================

\subsection{Mod\`ele logique de donn\'ees}
Le passage du mod\`ele objet au mod\`ele relationnel suit les r\`egles usuelles : une
classe devient une table, une composition devient une cl\'e \'etrang\`ere avec
suppression en cascade, un objet-valeur imbriqu\'e est aplati dans la table de son
propri\'etaire. Deux \'ecarts \`a ces r\`egles ont \'et\'e faits d\'elib\'er\'ement et sont
justifi\'es ci-apr\`es.

\begin{figure}[H]
\centering
\ajusterportrait{%
\begin{diagramme}
  \node (te) [umlcompact=3.9cm] at (0,0) {
    \umltitres{table}{events}
    \nodepart{two}
    \underline{event\_id} \,: TEXT \textsc{pk}\\
    fingerprint \,: TEXT \textsc{idx}\\
    incident\_id \,: TEXT \textsc{fk}\\
    occurred\_at, received\_at \,: TEXT\\
    source, category \,: TEXT\\
    severity, confidence \\
    payload \,: TEXT \tiny(JSON)
  };
  \node (ti) [umlcompact=4.1cm] at (5.4,0) {
    \umltitres{table}{incidents}
    \nodepart{two}
    \underline{incident\_id} \,: TEXT \textsc{pk}\\
    correlation\_key \,: TEXT \textsc{idx}\\
    status, severity \,: TEXT\\
    opened\_at, updated\_at \,: TEXT\\
    attack\_code, family \,: TEXT\\
    dangerousness \,: REAL\\
    risk\_score \,: REAL \textsc{idx}
  };
  \node (td) [umlcompact=4.1cm] at (11.0,0) {
    \umltitres{table}{decisions}
    \nodepart{two}
    \underline{decision\_id} \,: TEXT \textsc{pk}\\
    incident\_id \,: TEXT \textsc{fk}\\
    outcome \,: TEXT\\
    rationale \,: TEXT\\
    decided\_at \,: TEXT\\
    trace \,: TEXT \tiny(JSON)\\
    fallback \,: TEXT \tiny(JSON)
  };
  \node (ta) [umlcompact=4.1cm] at (11.0,-5.6) {
    \umltitres{table}{actions}
    \nodepart{two}
    \underline{action\_id} \,: TEXT \textsc{pk}\\
    incident\_id, decision\_id \,: \textsc{fk}\\
    verb, actuator, target \,: TEXT\\
    status \,: TEXT \textsc{idx}\\
    reversibility \,: TEXT\\
    rollback\_token \,: TEXT\\
    started\_at, finished\_at \,: TEXT
  };
  \node (tj) [umlcompact=4.3cm, umlaccentue] at (5.4,-5.6) {
    \umltitres{table, append-only}{audit\_log}
    \nodepart{two}
    \underline{seq} \,: INTEGER \textsc{pk}\\
    event\_type \,: TEXT \textsc{idx}\\
    incident\_id \,: TEXT \textsc{idx}\\
    actor \,: TEXT\\
    recorded\_at \,: TEXT\\
    payload \,: TEXT \tiny(JSON)\\
    prev\_hash, entry\_hash \,: TEXT
  };
  \node (tp) [umlcompact=3.7cm] at (0,-10.4) {
    \umltitres{table}{policies}
    \nodepart{two}
    \underline{policy\_id} \,: TEXT \textsc{pk}\\
    version \,: INTEGER\\
    source\_text \,: TEXT\\
    compiled \,: TEXT \tiny(JSON)\\
    fingerprint \,: TEXT\\
    compiled\_at \,: TEXT
  };
  \node (tn) [umlcompact=3.7cm] at (0,-5.6) {
    \umltitres{table}{notifications}
    \nodepart{two}
    \underline{notification\_id} \,: TEXT \textsc{pk}\\
    incident\_id \,: TEXT \textsc{fk}\\
    level, subject, body \,: TEXT\\
    created\_at \,: TEXT\\
    acknowledged\_at \,: TEXT \textsc{idx}
  };
  \node (tu) [umlcompact=3.7cm] at (5.4,-10.4) {
    \umltitres{table}{users}
    \nodepart{two}
    \underline{user\_id} \,: TEXT \textsc{pk}\\
    username \,: TEXT \textsc{uq}\\
    password\_hash \,: TEXT\\
    role, status, kind \,: TEXT\\
    poste, civility \,: TEXT
  };
  \node (ts) [umlcompact=3.7cm] at (11.0,-10.4) {
    \umltitres{table}{user\_sessions}
    \nodepart{two}
    \underline{session\_id} \,: TEXT \textsc{pk}\\
    user\_id \,: TEXT \textsc{fk}\\
    token\_fingerprint \,: TEXT\\
    expires\_at \,: TEXT \textsc{idx}
  };

  \draw[association] (te) -- node[mult, near start]{1..*} node[mult, near end]{1} (ti);
  \draw[association] (ti) -- node[mult, near start]{1} node[mult, near end]{0..*} (td);
  \draw[association] (ti.south) -- node[mult, near start]{1} node[mult, near end]{0..*} (tj.north);
  \draw[association] (td.south) -- node[mult, near start]{1} node[mult, near end]{0..*} (ta.north);
  \draw[association] (ti.south east) -- node[mult, near end]{0..*} (ta.north west);
  \draw[association] (ti.west) -- node[mult, near end]{0..*} (tn.north east);
  \draw[association] (tu) -- node[mult, near start]{1} node[mult, near end]{0..*} (ts);
\end{diagramme}}
\caption{Mod\`ele logique de donn\'ees}
\label{fig:mld}
\end{figure}

\paragraph{Premier \'ecart --- la charge utile d'origine est conserv\'ee.} La colonne
\texttt{events.payload} garde le message brut de la source, sans normalisation.
C'est une redondance apparente avec les colonnes normalis\'ees. Elle est n\'ecessaire
\`a deux usages : rejuger une d\'ecision \`a partir de ce que la source avait
r\'eellement dit, et rejouer un \'ev\'enement apr\`es correction d'un adaptateur.

\paragraph{Second \'ecart --- la trace de d\'ecision n'est pas \'eclat\'ee en tables.}
\texttt{decisions.trace} est stock\'ee en document structur\'e plut\^ot qu'en tables
normalis\'ees. La raison n'est pas la commodit\'e : une trace est un
\emph{t\'emoignage fig\'e}. L'\'eclater en tables la rendrait modifiable par jointure et
sensible \`a toute \'evolution ult\'erieure du sch\'ema, ce qui contredirait
l'opposabilit\'e recherch\'ee (ENF-02).

\subsection{Immutabilit\'e du journal d'audit}
Le journal ne se contente pas d'\^etre en ajout seul par convention : la contrainte
est port\'ee par le moteur de base de donn\'ees lui-m\^eme, au moyen de d\'eclencheurs qui
font \'echouer toute mise \`a jour et toute suppression. S'y ajoute un cha\^inage par
empreintes : chaque entr\'ee porte l'empreinte de la pr\'ec\'edente, de sorte que
l'alt\'eration d'une entr\'ee ancienne invalide toutes les suivantes.

\begin{figure}[H]
\centering
\begin{diagramme}
  \foreach \i/\x in {1/0, 2/4.2, 3/8.4} {
    \node (n\i) [umlcompact=3.4cm] at (\x,0) {
      \umltitre{entr\'ee $n_{\i}$}
      \nodepart{two}
      \tiny type, acteur, horodatage\\
      \tiny charge utile\\
      \texttt{prev\_hash}\\
      \texttt{entry\_hash}
    };
  }
  \node (nn) [umlsimple=1.6cm] at (12.0,0) {\umltitre{\ldots}};
  \draw[flot] (n1.east) -- node[role, above]{\texttt{entry\_hash}$_1$ $\rightarrow$ \texttt{prev\_hash}$_2$} (n2.west);
  \draw[flot] (n2.east) -- node[role, above]{\texttt{entry\_hash}$_2$ $\rightarrow$ \texttt{prev\_hash}$_3$} (n3.west);
  \draw[flot] (n3.east) -- (nn.west);
  \node[note=6.4cm, anchor=north] at (5.0,-2.2)
    {$\texttt{entry\_hash}_i = \mathrm{SHA\text{-}256}\big(\texttt{prev\_hash}_i \,\|\,
      \text{type} \,\|\, \text{acteur} \,\|\, \text{horodatage} \,\|\, \text{charge utile}\big)$\\[3pt]
     Modifier l'entr\'ee $i$ change $\texttt{entry\_hash}_i$, donc invalide
     $\texttt{prev\_hash}_{i+1}$, et de proche en proche toute la suite. La v\'erification
     de la cha\^ine est un parcours lin\'eaire, ex\'ecutable par un tiers.};
\end{diagramme}
\caption{Cha\^inage par empreintes du journal d'audit}
\label{fig:chainage}
\end{figure}

\subsection{Place du mod\`ele de langage}
\label{sec:assistant-conception}
Le syst\`eme comporte un assistant conversationnel destin\'e \`a l'analyste. Sa
conception ob\'eit \`a une r\`egle unique, pos\'ee avant toute autre consid\'eration
d'ergonomie : \textbf{le mod\`ele de langage ne produit aucun fait et ne choisit
aucune action}. Les faits sont \'etablis par le moteur d\'eterministe et lui sont
fournis ; il n'en assure que la mise en forme.

Cette s\'eparation est structurelle et non contractuelle : le composant qui \'etablit
les faits et celui qui les formule sont distincts, et le second n'a acc\`es \`a aucun
actionneur. Elle vaut \`a la fois pour la conformit\'e \`a l'exigence de non-invention
(EF-04) et pour la contrainte de souverainet\'e --- le syst\`eme reste enti\`erement
op\'erant si aucun mod\`ele n'est disponible, la formulation se repliant alors sur des
gabarits.

\begin{figure}[H]
\centering
\begin{diagramme}
  \node (q)  [action=2.6cm] at (0,0) {Question de l'analyste};
  \node (i)  [action=2.6cm] at (3.6,0) {Reconnaissance d'intention \tiny(d\'eterministe)};
  \node (f)  [action=2.9cm] at (7.4,0) {Collecte des faits \tiny(d\'ep\^ots, journal)};
  \node (m)  [action=2.9cm] at (11.4,0) {Mise en forme \tiny(gabarit ou mod\`ele local)};
  \node (r)  [action=2.6cm] at (15.2,0) {R\'eponse};
  \draw[flot] (q) -- (i); \draw[flot] (i) -- (f); \draw[flot] (f) -- (m); \draw[flot] (m) -- (r);
  \node[note=6.0cm, anchor=north] at (7.4,-1.4)
    {La fl\`eche entrante du composant de mise en forme ne transporte que des faits
     d\'ej\`a \'etablis. Aucune fl\`eche ne va de ce composant vers un actionneur : il ne
     peut, structurellement, d\'eclencher aucun geste.};
\end{diagramme}
\caption{Cha\^ine de traitement de l'assistant --- s\'eparation des faits et de leur formulation}
\label{fig:assistant}
\end{figure}

%%%==========================================================================
\section{Matrice de tra\c cabilit\'e}
%%%==========================================================================
La matrice du tableau~\ref{tab:tracabilite} relie chaque exigence \`a l'\'el\'ement de
conception qui la porte et \`a la figure o\`u il se lit. Elle sert deux usages : la
revue --- montrer qu'aucune exigence n'est rest\'ee sans r\'eponse --- et la
maintenance --- savoir ce qu'une modification met en jeu. Le tiret indique une
exigence port\'ee par la conception d'ensemble et non par un \'el\'ement unique.

\begin{table}[H]
\centering
\caption{Matrice de tra\c cabilit\'e --- exigence vers \'el\'ement de conception}
\label{tab:tracabilite}
\scriptsize
\begin{tabular}{@{}llp{5.4cm}l@{}}
\toprule
\textbf{Exigence} & \textbf{Vue} & \textbf{\'El\'ement de conception} & \textbf{Figure} \\
\midrule
EF-03 & Composants & Composant \emph{Enrichissement} & \ref{fig:composants} \\
EF-04 & Dynamique  & N\oe ud de d\'ecision \emph{contexte ancr\'e ?} & \ref{fig:activite} \\
EF-05 & Dynamique  & Activit\'e \emph{Planifier les actions} & \ref{fig:activite} \\
EF-06 & Statique   & \texttt{ActionSpec.parameters} & \ref{fig:classes-domaine} \\
EF-07 & Dynamique  & Message \emph{ex\'ecuter} sans message humain ant\'erieur & \ref{fig:sequence-systeme} \\
EF-08 & Composants & Composant \emph{UEBA} & \ref{fig:composants} \\
EF-09 & Composants & Composant \emph{UEBA} & \ref{fig:composants} \\
EF-10 & Statique   & \texttt{DecisionTrace.matchedConditions} & \ref{fig:classes-domaine} \\
EF-11 & Composants & Composant \emph{API et interface} & \ref{fig:composants} \\
EF-12 & Persistance& Table \texttt{audit\_log} & \ref{fig:mld} \\
EF-13 & Dynamique  & Message \emph{notification a posteriori} & \ref{fig:sequence-systeme} \\
EF-14 & Statique   & Invariant de construction d'\texttt{ActionSpec} & \ref{fig:classes-domaine} \\
EF-15 & Statique   & \texttt{PolicyVerdict}, table \texttt{policies} & \ref{fig:classes-domaine}, \ref{fig:mld} \\
EF-16 & Statique   & \texttt{Incident.riskScore()} & \ref{fig:classes-domaine} \\
EF-17 & Composants & Composant \emph{API et interface} & \ref{fig:composants} \\
EF-18 & Composants & Composant \emph{Adaptateur d'ingestion} & \ref{fig:composants} \\
EF-19 & Statique   & \texttt{DetectionEvent.fingerprint()} & \ref{fig:classes-domaine} \\
EF-20 & Statique   & \texttt{Incident.keyFor()}, \texttt{accepts()} & \ref{fig:classes-domaine} \\
EF-21 & Composants & Composant \emph{Sondes d'infrastructure} & \ref{fig:composants} \\
EF-22 & Composants & Composant \emph{Sondes d'infrastructure} & \ref{fig:composants} \\
EF-23 & Statique   & \texttt{AttackFamily.INFRASTRUCTURE} & \ref{fig:enumerations} \\
EF-25 & Dynamique  & Boucle de contr\^ole post-action & \ref{fig:sequence-boucle} \\
EF-26 & Dynamique  & N\oe ud \emph{coupe-circuit ferm\'e ?} & \ref{fig:activite} \\
EF-27 & Statique   & \texttt{FallbackPlanner}, \texttt{FallbackPlan} & \ref{fig:classes-repli} \\
EF-28 & Dynamique  & Activit\'e \emph{\'Emettre une alerte persistante} & \ref{fig:activite-repli} \\
EF-29 & Dynamique  & Activit\'e \emph{Ouvrir une demande de qualification} & \ref{fig:activite-repli} \\
\midrule
ENF-01 & D\'eploiement & N\oe ud unique, aucun lien sortant obligatoire & \ref{fig:deploiement} \\
ENF-02 & Persistance  & Cha\^inage par empreintes, d\'eclencheurs d'interdiction & \ref{fig:chainage} \\
ENF-03 & Statique     & \texttt{DecisionTrace} & \ref{fig:classes-domaine} \\
ENF-04 & Dynamique    & Absence de message humain bloquant & \ref{fig:sequence-systeme} \\
ENF-05 & Statique     & \texttt{SourceKind.REPLAY}, contrat d'idempotence & \ref{fig:enumerations}, \ref{fig:actionneurs} \\
ENF-06 & Statique     & Invariant global sur \texttt{Decision.actions} & tab.~\ref{tab:invariants} \\
ENF-07 & D\'eploiement & Topologie \`a n\oe ud unique & \ref{fig:deploiement} \\
ENF-08 & Paquetages   & Adaptateurs isol\'es dans \texttt{ingestion} & \ref{fig:paquetages} \\
ENF-09 & --- & Conception de l'interface d'exploitation & chap.~\ref{chap:implementation} \\
ENF-10 & Persistance  & Tables \texttt{users}, \texttt{user\_sessions} & \ref{fig:mld} \\
\bottomrule
\end{tabular}
\end{table}

%%%==========================================================================
\section{Conclusion du chapitre}
%%%==========================================================================
La mod\'elisation men\'ee dans ce chapitre fait appara\^itre une coh\'erence qui n'\'etait
pas acquise au d\'epart : \textbf{la contrainte de r\'eversibilit\'e se propage \`a tous
les niveaux du mod\`ele}. Elle est un attribut dans le diagramme de classes, un
invariant de construction dans les contrats, un n\oe ud de d\'ecision dans le
diagramme d'activit\'e, un \'etat dans la machine \`a \'etats de l'action, une colonne
dans le mod\`ele de donn\'ees, et une clause du contrat des actionneurs. Cette
propagation n'est pas une redondance : c'est ce qui rend la contrainte
incontournable par n'importe quel chemin d'ex\'ecution.

Deux limites de conception m\'eritent d'\^etre \'enonc\'ees ici plut\^ot que d\'ecouvertes \`a
la lecture des r\'esultats.

\begin{enumerate}
  \item \textbf{La mesure de r\'ef\'erence de la boucle de contr\^ole n'est pas
        persist\'ee.} Apr\`es un red\'emarrage, la boucle s'abstient sur les actions
        ant\'erieures. Le comportement est s\^ur --- elle pr\'ef\`ere ne rien faire plut\^ot
        que comparer \`a une r\'ef\'erence qu'elle n'a pas prise --- mais il r\'eduit la
        port\'ee r\'eelle d'EF-25 apr\`es incident d'exploitation.
  \item \textbf{La qualification d'une menace hors catalogue reste une
        proposition.} Le syst\`eme d\'eduit un nom, une famille et une criticit\'e des
        indicateurs, mais l'entr\'ee ne rejoint le catalogue qu'apr\`es validation
        humaine. Ce choix est d\'elib\'er\'e --- un catalogue qui s'enrichirait seul
        d\'eriverait sans contr\^ole --- mais il maintient une d\'ependance \`a la
        disponibilit\'e d'un analyste sur ce point pr\'ecis, et sur celui-l\`a seulement.
\end{enumerate}

Le chapitre suivant confronte ces mod\`eles \`a leur r\'ealisation.
